\section{Experimental Setup}
\label{sec:setup}

\yuji{Just updated a new version of the experimental setup, will fix more details in more rounds later}

Having constructed \textsc{RippleEval}, we use it to study how a controlled factual update affects connected factual knowledge in LLMs. Our experiments are designed around a single-edit setting: for each target fact, we independently update the model with a new to-be-update object and then evaluate how this update changes the model's behavior on related QA samples. This setup allows us to isolate the effect of one factual update while avoiding the confounding effects of large-scale continual training.

\subsection{Subject Models}

We evaluate four open-weight models spanning two families, two scales, and two instruction-tuning regimes:

\begin{itemize}
    \setlength\itemsep{0em}
    \item \textbf{Qwen3.5-2B} \cite{qwen3report}: a small base model from the Qwen3 series.
    \item \textbf{Qwen3.5-9B} \cite{qwen3report}: a mid-scale model from the same family, used to test how ripple effects change with model capacity within a fixed architecture.
    \item \textbf{Gemma-4-E4B-it} \cite{gemma4report}: a small instruction-tuned model from the Gemma-4 series, used to probe whether aggressive instruction-tuning reduces susceptibility to factual updates.
    \item \textbf{Gemma-4-31B-it} \cite{gemma4report}: a large instruction-tuned model from the same family, providing a higher-capacity counterpart to Gemma-4-E4B-it.
\end{itemize}

Using two families at two scales each allows us to test whether the observed ripple effects are consistent across model size and instruction-tuning, rather than specific to a single architecture.

\subsection{Updating Target Selection}

For each experiment, we select a target factual relation $(s,r,o)$ from the verified factual graph $\mathcal{G}_{fact}$. Each selected target is associated with the in-degree of its object entity $o$, which provides the factual connectivity score defined in Section~\ref{sec:dataset}. This score is used to organize target updates in later analyses.

The base model is reset before each target update. Therefore, each run measures the effect of a single factual update from the same pre-update model state, rather than the accumulated effect of multiple edits.


\subsection{Knowledge Update}

For each selected target fact $(s,r,o)$, we construct an update by replacing the original object $o$ with a plausible but structurally distinct target object $o'$. For example, a target fact such as \textit{CapitalOf(France, Paris)} may be updated to \textit{CapitalOf(France, Lyon)}. We use this substitution to create a controlled perturbation that is unlikely to be already stored as the model's factual knowledge. This allows us to isolate how an injected update affects connected factual knowledge, rather than conflating the effect with the model's prior memory of real-world facts or uncontrolled real-world logical associations.

To train the model on the target update, we generate QA-format instruction pairs that express the new relation $(s,r,o')$. For each target fact, we first generate 30 diverse QA templates and then apply paraphrase augmentation to expand them into 150 targeted update pairs. These pairs serve as the direct supervision for injecting the target object $o'$.

To reduce overfitting to the target update and preserve unrelated model behavior, the final update set contains 650 QA pairs:

\begin{itemize}
    \setlength\itemsep{0em}
    \item 150 targeted update pairs expressing the injected fact $(s,r,o')$;
    \item 400 neutral factual QA pairs sampled from distant, disconnected subgraphs in $\mathcal{G}_{fact}$;
    \item 100 out-of-domain irrelevant QA pairs, disjoint from the irrelevant evaluation set.
\end{itemize}

The generated QA prompts are fixed and reused before and after updating. This ensures that observed behavioral changes are not caused by prompt variation.

\subsection{Update Optimization}

We implement the factual update using Low-Rank Adaptation (LoRA) \cite{hu2021lora}. We use parameter-efficient tuning rather than full fine-tuning because our goal is to study how a small factual update propagates through the model's existing knowledge representations, rather than to overwrite the model globally. LoRA provides a natural gradient-based update while keeping most model parameters frozen, making it suitable for measuring ripple effects under limited parameter disturbance.

We train LoRA adapters on the update set by minimizing the cross-entropy loss of the injected object $o'$ under the target relation:

\[
\min_{\theta} \mathcal{L}_{CE}(P_{\theta}(o' \mid s,r)).
\]

Unless otherwise specified, we use LoRA rank $r=20$, scaling factor $\alpha=40$, dropout $p=0.1$, and apply LoRA to the \texttt{q\_proj} and \texttt{v\_proj} matrices. We optimize with AdamW \cite{loshchilov2019decoupled}, using learning rate $2.5 \times 10^{-4}$, batch size 4, and 8 epochs, stopping after convergence when the training loss falls below $10^{-3}$. Additional prompt templates and implementation details are provided in the Appendix.


\paragraph{Updating Success Evaluation: Injection Success Rate (ISR).}
Before analyzing unintended changes to neighboring facts, we first verify whether the update succeeds on the target fact itself. For an edit that replaces the original object $t$ with the injected target $t'$, ISR checks whether $t'$ receives the highest candidate score for the edited query $q_{h,r}$:
\begin{equation}
    \text{ISR}
    =
    \mathbb{I}
    \left[
    \arg\max_{y \in \mathcal{C}_{h,r}}
    S_\theta(y \mid q_{h,r})
    =
    t'
    \right],
    \label{eq:isr}
\end{equation}
where $\mathcal{C}_{h,r}$ is the candidate object set for the subject--relation query. We report the number of failed injections and compute downstream metrics both over all attempted edits and over the successful-edit subset, separating target-injection failures from unintended changes to related non-target facts.

\subsection{Preserving General Model Performance}

Because our goal is to study localized ripple effects without causing model collapse, we perform a global sanity check after each update. Specifically, we evaluate the edited model on a fixed set of 50 irrelevant factual questions that are disjoint from the target neighborhood and the update data. A valid update should preserve performance on this irrelevant set, indicating that the model has not undergone broad catastrophic forgetting.

This sanity check separates local ripple effects from global capability degradation. If the model fails broadly on unrelated questions, downstream errors cannot be confidently attributed to propagation through the factual graph. In our trials, accuracy on this irrelevant set remains stable, suggesting that the observed changes primarily reflect update-induced effects on connected factual knowledge rather than general model degradation.

\yuji{Previous version starts from here...}







\subsection{Evaluation Metrics: Measuring Error Propagation}
Standard accuracy metrics often fail to capture the systemic impact of an update. Therefore, we employ a three-tiered evaluation strategy to quantify the damage \cite{meng2022rome, cohen2023ripple}.


\paragraph{2. Error Propagation Rate (EPR).}
\weibing{To quantify the extent of error propagation rigorously, we define EPR using a paired-query protocol.}[todo: Replaced conditional probability notation with explicit paired queries] For each edited fact $(s,r,t)$, we construct a downstream query set $\mathcal{Q}_k(t)=\{(q_i, y_i^{fact})\}_{i=1}^{n_k}$, where $y_i^{fact}$ is the factual answer logically entailed by the original object $t$ along the directed N-to-1 path (or its downstream descendants).
We retain only the queries that the model answered correctly before the update: $\mathcal{B}_k=\{i: \text{Argmax}(P_{\theta_0}(\cdot | q_i))=y_i^{fact}\}$. \weibing{EPR then measures the fraction of these previously correct facts that are corrupted and fail to output the original factual answer after the update}[edited: Fixed incorrect EPR definition (C>W metric)]:
\begin{equation}
    \text{EPR}_k = \frac{|\{i\in \mathcal{B}_k: \text{Argmax}(P_{\theta'}(\cdot | q_i))\not= y_i^{fact}\}|}{|\mathcal{B}_k|}
    \label{eq:epr}
\end{equation}
\weibing{Here, a failure ($\not= y_i^{fact}$) indicates the model has forgotten or corrupted the correct factual logic, demonstrating a Correct-to-Wrong (C>W) flip. We additionally apply alias-normalized exact matching using Wikidata aliases to prevent formatting artifacts.}[todo: Corrected failure explanation to match C>W logic]
For rigorous paired auditing between Hub and Tail sources, we strictly control the evaluation distribution by uniformly sampling exactly $n=30$ neighbors per hop distance ($d \in [1, 5]$) for each update, ensuring balanced comparability across structural depths. \weibing{To rigorously track causality, we employ \textbf{Sub-tree Constraint Sampling}. Rather than sampling nodes independently per hop—which could break the logical causal chain—we ensure that if a node is evaluated at distance $d$, its exact causal parent must exist in the $d-1$ evaluation set. This strict path continuity guarantees that downstream errors are direct propagations from upstream corrupted hubs, allowing for conditional probability analysis ($P(E_{d} | E_{d-1})$).}[todo: Added sub-tree sampling constraint to defend against causal disconnect] A high EPR indicates that the model's internal consistency checks are failing, allowing errors to propagate to unrelated concepts.

\paragraph{3. Confidence Shift ($\Delta$ Conf).}
Finally, discrete label flips do not capture the full picture. A model might retain the correct answer but become uncertain, or alternatively, become highly confident in a new hallucination. To detect these ``Silent Failures'' \cite{guo2017calibration}, we track the probability mass assigned to the specific incorrect answer that the model generates post-update ($\hat{y}_{err}$):
\begin{equation}
    \Delta \text{Conf} = P_{post}(\hat{y}_{err}) - P_{pre}(\hat{y}_{err})
    \label{eq:delta_conf}
\end{equation}
\weibing{For open-vocabulary generation, $\hat{y}_{err}$ represents the normalized text string of the incorrectly inferred answer derived from the Generation Accuracy Protocol. To measure its confidence shift, we revert to the Candidate Scoring Protocol, computing the exact joint probability of $\hat{y}_{err}$ under both pre- and post-update models.}[todo: Aligned with the separated evaluation protocols] A positive shift ($\Delta \text{Conf} > 0$) implies the model is increasing its certainty in a specific hallucination.
To investigate \textit{why} popular knowledge might be inherently more susceptible to flipping (our first hypothesis), we examine the internal decision boundaries prior to any update. For a given target fact, we define the Logit Margin as the difference between the logit of the correct answer and the logit of the strongest incorrect competitor (runner-up):
\begin{equation}
    \text{Margin} = \text{Logit}(y_{correct}) - \text{Logit}(y_{runner\text{-}up})
    \label{eq:margin}
\end{equation}
To ensure a rigorous baseline and prevent confounding noise from poorly learned facts, all mechanistic evaluations (including Margin and $\Delta \text{Confidence}$) strictly pass through a baseline mask (denoted as \textit{Mask B}). This mask only retains samples where the clean pre-update model accurately predicts the correct target token at rank 1 (\texttt{clean\_accuracy == 1} and \texttt{clean\_correct\_token\_rank == 1}). A narrower initial margin indicates a more fragile representation within the feature space.

\paragraph{5. Attention Shift ($\Delta$ Attention Lift).}
To explain the mechanistic pathways of error propagation (our second hypothesis), we probe the model's internal attention matrices in the Llama-2 paired audit. For each evaluation prompt, we collect the generation attentions for the first generated token, extract the final transformer layer, and denote the resulting tensor by $A^{(L)} \in \mathbb{R}^{H \times Q \times K}$, where $H$ is the number of attention heads, $Q$ is the query-length axis for that generation step, and $K$ is the prompt context length. Let $S$ be the token span of the evaluated neighbor entity, identified by matching the tokenized string of the entity head in the input prompt. We first define the attention mass on $S$ as:
\begin{equation}
\mathrm{Mass}(S) = \sum_{k \in S} \frac{1}{HQ} \sum_{h=1}^{H} \sum_{q=1}^{Q} A^{(L)}_{h,q,k}
\label{eq:att_mass}
\end{equation}
We then normalize by the span-length baseline to obtain Attention Lift:
\begin{equation}
\mathrm{Lift}(S) = \frac{\mathrm{Mass}(S)}{|S|/K}
\label{eq:att_lift}
\end{equation}
and report the absolute post-update change
\begin{equation}
|\Delta \mathrm{Lift}(S)| = |\mathrm{Lift}_{post}(S) - \mathrm{Lift}_{pre}(S)|
\label{eq:delta_lift}
\end{equation}
For Llama-2, the probe uses cloze/completion prompts rather than chat-style QA prompts. This metric should therefore be interpreted as a mechanistic probe of last-layer, first-generation-step attention concentrated on the evaluated neighbor head-token span, not as a replacement for EPR.